\section{The gamma and sine tools used by the functional equation}
\label{sec:gamma-gateway}

We give the finite constructions needed by the next section, so that the
functional equation has no unmentioned dependence on a classical gamma
theory. These arguments also explain exactly which human-readable results
from the repository's \texttt{gamma\_ipra} package are reused~\cite{GammaRepo}.
They are not represented as existing kernel-checked analytic theorems.

\subsection{Reciprocal gamma from a convergent tail}
Let $\gamma$ have the name in \eqref{em:gamma-name}, and put
\[
 \ell(w)=\sum_{j=2}^{\infty}\frac{(-1)^{j+1}w^j}{j}
 \quad(|w|\le1/2).
\]
The tail after degree $d$ is at most
$|w|^{d+1}/((d+1)(1-|w|))$, and $|\ell(w)|\le |w|^2$.
The identity $e^{w+\ell(w)}=1+w$ follows from finite coefficient
comparison: the formal derivative equation is $(1+w)F'=F$, $F(0)=1$,
so $F_1=1$ and $(j+1)F_{j+1}=(1-j)F_j$. Every later coefficient is
zero. The geometric and factorial tails justify each coefficient
calculation at finite precision.

For an integer $R\ge1$, $|z|\le R$, and $K=2R$, define
\[
 P(z)=z\prod_{j=1}^{K}(1+z/j)
 \exp\left((\gamma-H_K)z+\sum_{j>K}\ell(z/j)\right).
 \tag{G1}\label{eq:recip-gamma}
\]
The last sum has tail bounded by $R^2/N$ after index $N\ge K$, because
$\sum_{j>N}j^{-2}\le1/N$, a telescoping finite rational inequality.
Increasing $K$ leaves (G1) unchanged, by the identity for $\ell$.
Define the finite expression
\[
 Q_N(z)=zN^{-z}\prod_{j=1}^N(1+z/j).
\]
With
\[
 M_R=R(1+R)^{2R}3^{2R^2+2R},\qquad D_R=R^2+R,
\]
the bounds $0\le\gamma\le1$ and $H_K\le K$ give
\[
 |P(z)|\le M_R,\qquad
 |Q_N(z)-P(z)|\le\frac{3M_RD_R}{N}
 \quad(N\ge\max(2R,D_R)).                      \tag{G2}\label{eq:gamma-product-rate}
\]
Indeed the exponent difference between $Q_N$ and (G1), after the first
$N$ factors have been aligned, has modulus at most
$R|H_N-\log N-\gamma|+R^2/N\le D_R/N$. The exponential difference
bound $|e^u-1|\le3|u|$ for $|u|\le1$ proves (G2).

If $|z+j|\ge\delta>0$ for $0\le j\le2R$, with $\delta\le1$, the
same product gives, for $N\ge2R$,
\[
 |P(z)|,|Q_N(z)|\ge
 \lambda_{R,\delta}:=
 \frac{\delta^{2R+1}}{(2R)!\,3^{2R^2+2R}}.
 \tag{G3}\label{eq:gamma-lower}
\]
Define $\Gamma(z)=1/P(z)$ on such certified charts. The approximants
\[
 \Gamma_N(z)=\frac{N!N^z}{z(z+1)\cdots(z+N)}
\]
then have error at most $3M_RD_R/(\lambda_{R,\delta}^2N)$.
Moreover $|\Gamma(z)|\ge1/M_R$: gamma is quantitatively nonzero there.

For local power series, rearrange the last sum in (G1) at a larger disk:
\[
 \sum_{j>K}\ell(z/j)
 =\sum_{r=2}^{\infty}\frac{(-1)^{r+1}z^r}{r}
                 \left(\zeta(r)-\sum_{j=1}^Kj^{-r}\right).
\]
On $|z|\le R$, $K=2R$, the tail after degree $d$ is bounded by
$(R^2/K)2^{1-d}$. Finite convolutions with the exponential series give
an entire name $P$ with local coefficient majorants. In particular,
the finite prefactor has coefficient absolute sum at radius $R$ at
most $R(1+R)^{2R}$, and the exponent's corresponding sum is at most
$2R^2+2R$. Finite convolution and the exponential majorant therefore
bound the coefficient absolute sum directly by $M_R$; this is not
inferred from a supremum norm. Constructing on $|z|\le2R$ gives the
coefficient bound $M_{2R}$ and its degree-$d$ error on $|z|\le R$ by
$M_{2R}2^{-d-1}$. Geometric inversion using (G3) gives local power
series for gamma wherever it is defined.

Finite cancellation gives
\[
 zQ_N(z+1)=\left(1+\frac{z+1}{N}\right)Q_N(z).
\]
Together with $Q_N(1)=(N+1)/N$, this proves
\[
 P(z)=zP(z+1),\qquad P(1)=1,
 \quad\Gamma(z+1)=z\Gamma(z),\quad\Gamma(n+1)=n!.
 \tag{G4}\label{eq:gamma-recurrence}
\]
The entire recurrence also gives
$P(z)=z(z+1)\cdots(z+n)P(z+n+1)$ and
$P'(-n)=(-1)^n n!$. Thus gamma's poles at the nonpositive integers are
simple, with residues $(-1)^n/n!$; this is a statement about a punctured
local expression, not a value assigned at the pole.

\subsection{Euler's integral from finite beta integrals}
Suppose $\Re z\ge1/A$ and $|z|\le R$, with positive integers $A,R$.
Finite integration by parts gives
\[
 \int_0^1t^{z-1}(1-t)^N\,dt
       =\frac{N!}{z(z+1)\cdots(z+N)}.
 \tag{G5}\label{eq:finite-beta}
\]
To verify this before taking endpoints, integrate on $[\epsilon,1]$.
An antiderivative of $t^{z-1}$ is $t^z/z$, and the lower endpoint has
bound $A\epsilon^{1/A}$. Each integration by parts reduces $N$ by one
and increases $z$ by one. The finite recurrence ends at $N=0$.
The error for all these steps is assigned before induction on $N$.

Multiply (G5) by $N^z$ and substitute $u=Nt$. The result is the integral of
$u^{z-1}(1-u/N)^N$ on $[0,N]$. For $0\le u\le N/2$, the logarithm
series gives
\[
 0\le -N\log(1-u/N)-u\le u^2/N,
 \quad
 0\le e^{-u}-(1-u/N)^N\le e^{-u}u^2/N.
\]
For the rest of $[0,N]$, both kernels lie between zero and $e^{-u}$.
Also
\[
 \int_0^\epsilon |u^{z-1}e^{-u}|\,du\le A\epsilon^{1/A},\qquad
 \int_T^\infty |u^{z-1}e^{-u}|\,du
       \le2^{R+1}R!e^{-T/2}.                   \tag{G6}\label{eq:gamma-integral-tails}
\]
For the second inequality use $u^R e^{-u}\le2^RR!e^{-u/2}$ on
$u\ge1$, obtained from one term of the exponential series. Splitting the
kernel comparison at $1$ and $N/2$ now yields
\[
 \left|\Gamma_N(z)-\int_0^\infty u^{z-1}e^{-u}\,du\right|
 \le\frac{1+2^{R+3}(R+2)!}{N}+2^{R+1}R!e^{-N/4}.
\]
For this comparison take $N\ge2$. Together with (G2)--(G3), this proves the Euler integral representation.
Every integral here means a finite sum with the tails (G6); on the
remaining compact interval its derivative supplies a rational Lipschitz
bound. Differentiation with respect to $z$ introduces $(\log u)^j$.
Near zero use
$|\log u|^j\le j!(2A)^j u^{-1/(2A)}$, and at infinity use
$(\log u)^j\le u^j$. These give explicit PR derivative and Taylor
bounds by the same argument.
For real $x>0$ the integral is positive; the interval $[1,2]$ alone
gives $\Gamma(x)\ge1/18$.

\subsection{Reflection and duplication with finite product estimates}
For clarity we include the constructive sine-product step. For
$m=2J+1$ the multiple-angle recurrence gives a degree-$J$ polynomial in
$\sin^2u$ for $\sin(mu)/(m\sin u)$. Its $J$ known roots are
$\sin^2(k\pi/m)$, $1\le k\le J$, and its constant term is one.
They are apart: on this interval the complementary-angle chord bound
\eqref{eq:sine-chord} gives $\cos v\ge1/m$, so consecutive sine
values differ by at least $\pi/m^2$. Polynomial division by these
known roots proves the finite identity
\[
 \operatorname{sinc}(mu)=\operatorname{sinc}(u)
       \prod_{k=1}^{J}\left(1-\frac{\sin^2u}{\sin^2(k\pi/m)}\right),
 \qquad \operatorname{sinc}(u)=\sum_{j\ge0}\frac{(-1)^ju^{2j}}{(2j+1)!}.
 \tag{G7}\label{eq:finite-sine-product}
\]
Multiplying through by the indicated sinc avoids division by $\sin u$.

Here is a sufficient quantitative passage to Euler's product. For
$|z|\le R$, set $m=8K^2R^2+1$, $u=\pi z/m$, $v_k=\pi k/m$.
By the chord bound and factorial series,
\[
 a_k:=\frac{\sin^2u}{\sin^2v_k},\quad
 |a_k|\le16R^2/k^2,\quad
 \left|a_k-z^2/k^2\right|
 \le\frac{96R^2}{m^2}\left(\frac{R^2}{k^2}+1\right)\quad(k\le K).
\]
For the second estimate write sine as its argument times sinc; on
$|w|\le1$, $|\operatorname{sinc}w-1|\le |w|^2/2$, and
$\operatorname{sinc}v_k>1/2$. A common denominator and
$\pi^2<16$ give the displayed constant. The omitted factors have
$\sum_{k>K}|a_k|\le16R^2/K$.
Replacing factors one at a time bounds the resulting product difference
by $e^{34R^2}$ times the sum of these discrepancies. The prefactor sinc
in (G7) contributes at most $8R^2e^{32R^2}/m^2$.
Thus, for example, the generous rational bound
\[
 \left|\operatorname{sinc}(\pi z)
        -\prod_{k=1}^K(1-z^2/k^2)\right|
 \le\frac{2^{100(R+1)^2}}{K}                    \tag{G8}\label{eq:sine-rate}
\]
is valid. This is a PR product modulus proved from a finite polynomial
identity.

Exact cancellation in the Euler gamma approximants gives
\[
 \Gamma_N(z)\Gamma_N(-z)
       =-\frac1{z^2\prod_{k=1}^N(1-z^2/k^2)}.
\]
Use (G8), (G2)--(G4) to obtain
\[
 \Gamma(z)\Gamma(1-z)=\frac{\pi}{\sin\pi z}.
 \tag{G9}\label{eq:gamma-reflection}
\]
Division is licensed by a supplied separation $\delta$ from the finitely
many integers in the relevant bounded disk. One explicit bound, for an
integer $L\ge2R$, is
\[
 |\sin\pi z|\ge
       \frac{\delta^{2L+1}2^{-4R^2}}{(L!)^2}.
 \tag{G10}\label{eq:sine-apartness}
\]
Indeed the first $L$ factors of (G8) have product at least
$\delta^{2L}/(L!)^2$; for the rest use
$\log(1-t)\ge-2t$ on $0\le t\le1/4$ and
$\sum_{k>L}R^2/k^2\le R^2/L$. Multiply by $|\pi z|\ge\delta$.
At $z=1/2$, positivity of the Euler integral selects
$\Gamma(1/2)=\sqrt\pi$, where the positive square root is computed
by bounded rational bisection.

Duplication needs no asymptotic normalization. Let
$U_N(z)=\Gamma_N(z)\Gamma_N(z+1/2)/\Gamma_{2N+1}(2z)$.
Grouping its even and odd denominator factors gives exactly
\[
 \frac{U_N(z)}{U_N(1/2)}
       =\left(\frac{N}{2N+1}\right)^{2z-1}.
\]
Since $|\log(2N/(2N+1))|\le1/N$, the right side converges to
$2^{1-2z}$ with a PR modulus on each bounded disk. The gamma convergence
bounds and $\Gamma(1/2)=\sqrt\pi$ give
\[
 \Gamma(z)\Gamma(z+1/2)=2^{1-2z}\sqrt\pi\,\Gamma(2z).
 \tag{G11}\label{eq:gamma-duplication}
\]
First this is proved on charts separated from all displayed poles;
cleared local expressions give its meromorphic reading elsewhere.

\subsection{A zeta generating series and the derivative at one}
For $|z|\le r<1$, expansion of (G1) after the recurrence (G4) gives
\[
 \log\Gamma(1+z)
   =-\gamma z+\sum_{k=2}^{\infty}
                    \frac{(-1)^k\zeta(k)}{k}z^k.
 \tag{G12}\label{eq:log-gamma-zeta}
\]
Here the left side is the particular analytic logarithm with value zero
at $z=0$: exponentiating the right side equals $\Gamma(1+z)$.
For a proof, expand each local $\log(1+z/n)$ and rearrange finite
rectangles; the absolute double tail is bounded by
$\sum_{k>d}2r^k/k\le2r^{d+1}/((d+1)(1-r))$ because $\zeta(k)\le2$
for $k\ge2$. The remaining finite $n$-tails are bounded by the Dirichlet
estimates. Thus no rearrangement without a modulus occurs.
This proves $\Gamma'(1)=-\gamma$ and, on $|z|<1$,
\[
 \frac{\Gamma'(1+z)}{\Gamma(1+z)}
   =-\gamma+\sum_{k=2}^{\infty}(-1)^k\zeta(k)z^{k-1},
 \qquad
 \left(\frac{d}{dz}\right)^j\!\log\Gamma(1+z)\bigg|_{z=0}
 =(-1)^j(j-1)!\zeta(j)\quad(j\ge2).
\]
For derivatives, shrink the displayed radius from $r$ to any certified
smaller radius and bound $k^j(r'/r)^k$ by a factorial-exponential
estimate. This supplies a PR modulus for every fixed derivative order.
